\section{Background and Problem Formulation}
\label{sec:background}
In this section, we review the three categories of join discovery from Table~\ref{tab:intro-landscape}: non-LLM, LLM-based, and hybrid.
For each category, we describe how it decides joins and where it falls short on enterprise databases.
We then review how profiling metadata is collected and provided to LLMs, and formalize cost-bounded join key discovery.

\subsection{Non-LLM Join Discovery}
\label{sec:bg_nonllm}
Non-LLM methods decide joins from data statistics and schema metadata.
They fall into two groups: key and inclusion-dependency discovery, which decides joins inside one database, and joinable-table search, which ranks candidate columns for a query column in a data lake.

\introparagraph{Key and Inclusion-Dependency Discovery} A unary \emph{inclusion dependency} (IND) $A \subseteq B$ holds if every value of column $A$ also appears in column $B$.
If $B$ is also unique, then $A$ is a foreign-key candidate that references $B$.
IND discovery algorithms check this condition for all column pairs, and scale to databases with very many tables by sorting and sharing value sets across checks~\cite{Tschirschnitz17}.
Real data rarely satisfies exact INDs, so approximate INDs relax the condition to a containment threshold.
Every IND is only a candidate, and most candidates are not foreign keys.
Learned detectors therefore rank candidates with features, such as the uniqueness of $B$, the containment of $A$ in $B$, the similarity of column names, and the distribution of values~\cite{Rostin09,ZhangHOPS10}.
HoPF~\cite{HoPF} detects primary and foreign keys together, so that a column chosen as a primary key constrains which foreign keys may reference it.
Auto-BI~\cite{AutoBI} goes one step further.
It first trains a local classifier that predicts, for each column pair, whether it is an N:1 join.
It then selects a global join graph that maximizes the combined classifier scores under structural constraints, a variant of the minimum-cost arborescence problem.
The global step removes many locally plausible but inconsistent joins.

\introparagraph{Joinable-Table Search} Data lake systems treat each column as a \emph{domain}, the set of its distinct values, and rank the columns of a lake by their overlap with a query column.
LSH Ensemble~\cite{DBLP:journals/pvldb/ZhuNPM16} uses set containment $|Q \cap X| / |Q|$ as the overlap measure.
Unlike Jaccard similarity, containment does not penalize a large domain $X$ for values outside $Q$, which matters because domain sizes follow a power-law distribution.
LSH Ensemble indexes MinHash sketches of all domains in partitions of similar size, and answers containment queries over hundreds of millions of domains in seconds.
JOSIE~\cite{DBLP:conf/sigmod/ZhuDNM19} returns the exact top-$k$ columns by set overlap and chooses between reading posting lists and reading candidate sets with a cost model.
MATE~\cite{MATE} extends the search to multi-column keys.
DeepJoin~\cite{2023_dong_deepjoin} embeds columns with a pre-trained language model and searches them with an approximate nearest-neighbor index, which also finds joins with similar but not identical values.
Freyja~\cite{Freyja} learns a join quality metric from column profiles instead of raw values.

\introparagraph{Limitations} Non-LLM methods are cheap, but they decide joins from syntactic evidence, and enterprise databases defeat this evidence in four ways.
First, accidental overlap creates false joins.
Surrogate keys that count from 1, status codes, years, and flags are contained in each other, and a unique surrogate column looks like a valid target for all of them.
LakeBench~\cite{LakeBench} shows the same problem for joinable-table search: a column of street names overlaps a query column of last names, and JOSIE ranks it among the top results although the two columns are not joinable.
Second, real joins look imperfect.
Staging tables, snapshots, and deleted rows make containment partial, and formatting differences make exact overlap zero.
On LakeBench queries whose joins have only fuzzy overlap, JOSIE reaches 0\% top-10 accuracy and DeepJoin reaches 33\%~\cite{LakeBench}.
LakeBench concludes that no method is effective in all cases, because joinable columns must both mean the same thing and overlap in value~\cite{LakeBench}.
Third, names do not resolve meaning.
One key appears under different names, and different keys share generic names such as \texttt{ID}, \texttt{CODE}, or \texttt{DATA}.
Column descriptions and business context would resolve these cases, but these methods do not read them.
Fourth, learned models need labels.
The detectors and Auto-BI learn from labeled joins, and labels for a new enterprise schema are expensive: LakeBench spent more than 7,500 human hours to label its queries~\cite{LakeBench}.
Joinable-table search, in turn, returns a ranked list per query column, not a join decision with a key direction.

\subsection{LLM-based Join Discovery}
\label{sec:bg_llm}
LLMs decide what columns mean.
Foundation models match entities and schema elements across tables without task-specific training~\cite{dataWrangleNarayan2022}, and they read names, descriptions, and sample values together.
This resolves the cases that defeat non-LLM methods: an LLM recognizes that \texttt{ENGAGEMENT} holds project codes, that a surrogate \texttt{ID} does not reference a status table, and that a street name is not a last name.

\introparagraph{All-Pairs Prompting} Current practice in enterprise data platforms, including the production pipeline we study, discovers joins in three LLM steps.
First, the LLM annotates every table: it tags its columns with semantic types and writes a description.
Second, it tags each table with the \emph{concepts} of a business data model, such as project or invoice, whose data the table holds.
Third, it discovers joins pair by pair.
For each pair of tables, one prompt removes the columns that cannot be keys, a second prompt selects the candidate key columns of one table, and one further prompt per candidate lists all columns of the other table with their descriptions and sample values and returns the matching columns.
Warehouse queries then score the matches by value overlap and uniqueness and assemble a composite key greedily.
The first two steps make one LLM call per table or per batch of tables, so their cost is linear in the number of tables.
The third step dominates the cost.

\introparagraph{Limitations} All-pairs prompting is accurate, but it faces four limitations.
First, its cost grows quadratically.
A database with $n$ tables needs $n(n-1)/2$ pair evaluations, about 32M for 8K tables, and each evaluation makes $2 + c$ LLM calls for $c$ candidate columns, most of which carry a full schema and sample values.
In our databases, only \todo{Z}\% of the pairs join, so almost all calls return ``no join''.
Second, its calls are redundant.
If $k$ tables reference one key, the LLM answers $O(k^2)$ questions whose answers follow from $k$ references to the owner of the key.
Third, its answers are independent, so nothing makes them consistent.
The LLM may confirm that $A$ joins $B$ and $B$ joins $C$ on the same key but reject that $A$ joins $C$.
It also cannot see global evidence that is obvious across pairs: in one of our databases, it linked a generic column named \texttt{DATA} to \todo{1,826} other columns, one pair at a time.
The resulting join graph contains cliques of redundant edges and false edges that are hard to remove afterwards.
Fourth, it is incremental only in the trivial sense.
It skips pairs it has already evaluated, so a new table costs $n$ new pair evaluations, and every evaluation re-derives facts that earlier pairs already established, such as which column of a table holds customer identifiers.

\subsection{Hybrid Join Discovery}
\label{sec:bg_hybrid}
Hybrid methods combine cheap methods with the LLM to cut the number of LLM calls.
They either narrow the search space before the LLM, or route each candidate to a cheap proxy or to the LLM.

\introparagraph{Narrowing the Search Space} Schuyler~\cite{Schuyler} clusters the tables of a relational database into topics.
It fine-tunes a language model with a triplet loss on self-supervised signals from the schema, embeds each table, and clusters the embeddings, without labeled data.
Its authors note that such clusters reduce the search for joinable and unionable tables, because integration can match groups before matching tables.
The concept step of all-pairs prompting follows the same idea.
Magneto~\cite{Magneto} narrows schema matching between two tables: a small language model retrieves candidate column matches, and an LLM reranks them.

\introparagraph{Routing Candidates} LLM-powered data processing systems, such as LOTUS~\cite{LOTUS} and Palimpzest~\cite{Palimpzest}, provide a \emph{semantic join}: a join of two tables whose condition is a natural-language predicate evaluated by an LLM on each pair of records.
To avoid evaluating the full cross product, they use a \emph{model cascade}~\cite{FrugalGPT}.
A cheap proxy, typically the cosine similarity of record embeddings, scores each pair.
Pairs with a low score are rejected, pairs with a high score are accepted, and only pairs in between go to the LLM.
SUPG~\cite{SUPG} and BARGAIN~\cite{BARGAIN} set the two thresholds from a sample of LLM labels, so that the output reaches a target precision or recall of the all-LLM result with probability at least $1-\delta$.
The savings depend on the proxy.
When embeddings do not encode the join condition, most pairs fall between the thresholds: on real-world documents, BARGAIN still defers about 80\% of pairs to the LLM~\cite{FDJ}.
Featurized-Decomposition Join (FDJ)~\cite{FDJ} replaces the single proxy with a \emph{featurized decomposition}.
An LLM extracts features from each record in one linear pass, such as dates, locations, and names.
A logical expression in conjunctive normal form over cheap comparisons of these features then prunes non-matching pairs, with thresholds set to keep statistical guarantees.
FDJ reduces cost by up to $10\times$ over BARGAIN~\cite{FDJ}.

\introparagraph{Limitations} Hybrids reduce the cost per pair, but they fall short for join discovery in two ways.
First, their cheap stage inherits the failures of non-LLM methods.
Topical clusters are not key domains: a customer or project key is referenced from every topic, so a search restricted to clusters loses joins, and a search across clusters returns to a quadratic number of pairs.
A proxy based on overlap or embeddings passes accidental overlaps to the LLM and drops joins whose values do not overlap before the LLM sees them.
Second, semantic join operators answer a different question.
They evaluate a \emph{given} predicate over the records of \emph{two given} tables.
Join discovery must find which of thousands of tables join and on which columns.
Applied to table or column pairs, a semantic join still judges every pair independently.
It does not use the structure of join keys: the referenced side is unique, the referencing side is contained in it, and many tables share one key.
FDJ shows that linear per-record extraction can replace quadratic per-pair decisions.
\ourmethod applies this principle to schemas, with key domains as the extracted structure.

\subsection{Profiling Metadata for LLMs}
\label{sec:bg_metadata}
All methods above consume column profiles, and LLM-based methods must also fit a description of each column into a prompt.
We review how dependencies are discovered and stored, how LLMs annotate columns, and how metadata is selected for a prompt.

\introparagraph{Dependency Discovery and Profile Stores} Data profiling computes metadata such as unique column combinations (UCCs), functional dependencies (FDs), and INDs~\cite{ProfilingSurvey}.
Metanome~\cite{Metanome} bundles discovery algorithms for these dependencies.
HyFD~\cite{HyFD} and HyUCC~\cite{HyUCC} combine approximation on samples with validation to find all minimal FDs and UCCs, and BINDER~\cite{BINDER} finds all INDs by partitioning the value sets of all columns.
The number of candidate dependencies grows exponentially with the number of columns, so exact discovery over a full warehouse is expensive.
Approximate algorithms reduce this cost.
FAIDA~\cite{FAIDA} tests IND candidates on HyperLogLog sketches and samples, and reports no false positives in its evaluation; Pyro~\cite{Pyro} discovers approximate FDs and UCCs that tolerate errors in the data.
DynFD~\cite{DynFD} maintains FDs as the data changes, instead of discovering them again.
Metacrate~\cite{Metacrate} stores them as a queryable database of data profiles.
Wan and Han~\cite{ValueDriven} argue that a valid dependency is not always a useful one: discovery should target the dependencies whose value for a downstream task exceeds the cost to discover and maintain them.

\introparagraph{LLM-based Profiling and Annotation} LLMs annotate columns with semantic types without task-specific training~\cite{DBLP:conf/vldb/KoriniB23,ArcheType}, and write column descriptions for catalogs and downstream tasks~\cite{TACO}.
Dependencies can also come from the LLM, not the data: LLMClean~\cite{LLMClean} lets an LLM derive ontological FDs from a table and uses them to detect errors.
DBAutoDoc~\cite{DBAutoDoc} documents databases without declared keys.
It combines statistical analysis with iterative LLM refinement, which propagates corrections along the schema's dependency graph, and its statistical stage improves foreign-key detection by 23 F1 points over an LLM-only approach.

\introparagraph{Selecting Metadata for LLM Prompts} A profile helps an LLM only if its relevant part fits into the prompt.
Shkapenyuk et al.~\cite{AutoMetadataT2S} summarize database profiles with an LLM and select schema elements for text-to-SQL with a recall-oriented schema linking algorithm.
On BIRD with GPT-4o, profile-based metadata raises execution accuracy from 49.8\% to 61.2\%, and schema linking raises it to 63.2\%~\cite{AutoMetadataT2S}.
REVEAL~\cite{REVEAL} annotates a column with a compact subset of the other columns of its table, selected for relevance and diversity, instead of the whole table, and improves column type and column property annotation.

\introparagraph{Limitations} These works either collect dependencies or select context, but not both for one task.
Profiling systems discover all valid dependencies of a type and leave their use to the application.
Context selection methods choose columns or schema elements, not dependencies.
To our knowledge, no method decides which dependencies to discover, and which to show the LLM, for a given column and task under a token budget.
Join discovery needs exactly this decision.
An IND into a unique column suggests a reference, the description of the referenced key tells the LLM what the referencing column means, and most other dependencies do not affect the decision.
\ourmethod makes this selection with key domains: a ring prompt contains the owner of a key domain and its candidate members, together with the profile evidence that connects them.

\subsection{Problem Formulation}
\label{sec:problem}
Let $\mathcal{D}$ be the database of one business, with tables $\mathcal{T} = \{T_1, \ldots, T_n\}$, and let $\mathcal{C}$ be the set of all their columns.
A business typically has 50 to 8K tables.
Businesses are independent: every definition below is relative to one $\mathcal{D}$, and joins across businesses are out of scope.
Each column $c \in \mathcal{C}$ has a name, an optional LLM annotation (semantic type and description), and a set of values $V(c)$.
We write $u(c) = |\mathrm{distinct}(V(c))| / |V(c)|$ for its uniqueness, and $\mathrm{cont}(a, b) = |V(a) \cap V(b)| / |V(a)|$ for the containment of $a$ in $b$.

\begin{definition}[Join Key]
A \emph{join key} $a \rightarrow b$ is an ordered pair of columns $a, b \in \mathcal{C}$ of different tables, where $b$ identifies a business entity and $a$ refers to the same entity.
We call $b$ the \emph{referenced} column and $a$ the \emph{referencing} column.
\end{definition}

The definition is semantic on purpose.
Uniqueness and containment are evidence for a join key, but neither is necessary nor sufficient on enterprise data, as Section~\ref{sec:bg_nonllm} shows.
We therefore define the ground truth by an oracle.

\begin{definition}[LLM Oracle]
An \emph{oracle} $\mathcal{O}$ is an LLM with a fixed prompt that, given two tables, returns the set of join keys between them.
The oracle join set $J^* = \bigcup_{i<j} \mathcal{O}(T_i, T_j)$ is the result of all-pairs prompting.
Computing $J^*$ costs $n(n-1)/2$ oracle calls.
\end{definition}

Using the oracle as the reference makes the goal precise: a cheaper method should return what all-pairs prompting would return, at a fraction of its cost.
It also matches how such systems are evaluated in practice, where the all-LLM result is the reference for cascades~\cite{SUPG,BARGAIN,FDJ}.
We additionally report precision and recall against human-labeled joins on a sample (Section~\ref{sec:experiments}).

\begin{definition}[Key Domain and Ring]
A \emph{key domain} is a business entity that columns refer to, such as a customer or a project.
A \emph{ring} $R = (o_R, M_R)$ represents one key domain by an \emph{owner} column $o_R \in \mathcal{C}$ that identifies the entity, and a set of \emph{member} columns $M_R \subseteq \mathcal{C} \setminus \{o_R\}$ that refer to it.
A ring induces the join keys $J(R) = \{ m \rightarrow o_R \mid m \in M_R \}$.
A set of rings $\mathcal{R}$ induces the join graph $G(\mathcal{R})$ with edges $\bigcup_{R \in \mathcal{R}} J(R)$.
\end{definition}

Two members of the same ring also join with each other, through their shared owner.
$G(\mathcal{R})$ does not store these joins as edges, because they follow by transitivity.
So a key referenced by $k$ columns produces $k$ edges in $G(\mathcal{R})$, a star around its owner, instead of the $O(k^2)$ edges of a clique.

\introparagraph{Cost Model} Let $\mathrm{cost}(P)$ be the cost of an LLM prompt $P$, measured in tokens and converted to dollars by the price of the model.
A method $\mathcal{A}$ issues a sequence of prompts $P_1, \ldots, P_m$ and returns a join set $\hat{J}$.
Its LLM cost is $\mathrm{cost}(\mathcal{A}) = \sum_{i=1}^{m} \mathrm{cost}(P_i)$.
We also report its number of LLM calls $m$ and its end-to-end time, including profiling and all non-LLM steps.
The recall of $\hat{J}$ with respect to the oracle is $\mathrm{rec}(\hat{J}) = |\hat{J} \cap J^*| / |J^*|$, where a member-to-member join implied by a ring counts as found.

\introparagraph{Optimization Objective} Given $\mathcal{D}$, an oracle $\mathcal{O}$, a recall target $\gamma \in (0, 1]$, and a confidence level $1-\delta$, find a set of rings $\mathcal{R}$, and thereby a join set $\hat{J} = G(\mathcal{R})$, that minimizes the LLM cost:
\begin{equation}
\min_{\mathcal{A}} \ \mathrm{cost}(\mathcal{A}) \quad \text{s.t.} \quad \Pr\left[\mathrm{rec}(\hat{J}) \geq \gamma\right] \geq 1 - \delta,
\end{equation}
subject to five requirements.
\begin{itemize}[leftmargin=*]
\item \emph{(R1) Recall-guaranteed:} $\hat{J}$ keeps at least a fraction $\gamma$ of the oracle joins with probability at least $1-\delta$, where the probability is over the random samples the method draws.
\item \emph{(R2) Domain-proportional:} the number of LLM calls grows with the number of key domains and uncertain candidates, not with the number of table pairs.
\item \emph{(R3) Label-free:} the method needs no labeled joins and no corpus of similar schemas. It may learn only from the oracle answers it pays for.
\item \emph{(R4) Consistent:} $\hat{J}$ is induced by rings, so every referencing column points to the owner of its key domain, and joins implied by transitivity are not stored as separate edges.
\item \emph{(R5) Online:} when a table arrives, the method updates $\hat{J}$ at a cost that grows with the size of the new table and of the rings it affects, not with $n$, and the final $\hat{J}$ does not depend on the order in which tables arrive.
\end{itemize}
A precision constraint can be added in the same form; we report precision for all methods.

\introparagraph{Online Setting} Tables arrive over time.
Let $T_1, T_2, \ldots$ be the order of arrival, $\mathcal{D}_t = \{T_1, \ldots, T_t\}$ the business after $t$ arrivals, and $J^*_t$ the oracle join set of $\mathcal{D}_t$.
On the arrival of $T_{t+1}$, a method receives its columns and their profiles and annotations, and returns the updated join set $\hat{J}_{t+1}$.
It may add edges incident to $T_{t+1}$, and it may also revise earlier edges, for example when $T_{t+1}$ holds a better owner for an existing ring.
Requirement R1 must hold after every arrival, with $J^*_t$ as the reference, and the objective is the total LLM cost over the whole sequence.
The batch problem above is the special case of a single arrival that contains all tables.

In Table~\ref{tab:intro-landscape}, non-LLM methods meet R2 and R4 at best but have no oracle-level semantics, and Auto-BI violates R3.
All-pairs prompting is the oracle itself and violates R2, R4, and R5, because each arrival costs $n$ new pair evaluations.
Hybrid semantic joins meet R1 and R3 for a given pair of tables, but applied to join discovery they violate R2, R4, and R5.
Joinable-table search meets the cost part of R5, because a new column is a single index lookup, but it returns ranked candidates rather than joins.
